\subsection{Input handler}

\subsubsection{Purpose}
The main purpose of the input handler is to run actions with the inputs of the client.

\subsubsection{Functions}

% Optional overview paragraph
This module provides the following functions:

\begin{itemize}
    \item \textbf{Init(parameters)}  
    \textit{Description:} Briefly describe what this function does.  
    \textit{Inputs:} List parameters with types and purpose.  
    \textit{Outputs:} Return type and description.  
    \textit{Exceptions / Errors:} Any errors this function may raise.  
    \textit{Usage Example:} Optional small snippet showing a typical call.
    
    \item \textbf{FunctionName2(parameters)}  
    \textit{Description:} ...  
    \textit{Inputs:} ...  
    \textit{Outputs:} ...  
    \textit{Exceptions / Errors:} ...  
    \textit{Usage Example:} ...
\end{itemize}



\begin{lstlisting}[language=Lua, caption={Example: handleInput function}]
function handleInput(input)
    if input == "jump" then
        character:Jump()
    end
end
\end{lstlisting}

\subsubsection{Diagram}

\begin{tikzpicture}

% Nodes
\node (start) [startstop] {Initialize \& Detect Inputs};
\node (detect) [decision, below of=start] {Input Detected?};
\node (processEvent) [decision, below of=detect] {Game Processed Event?};
\node (handleEvent) [process, right of=processEvent, xshift=4cm] {Handle Game Event};
\node (checkList) [process, below of=processEvent] {Check Bind List};
\node (fireFunc) [process, below of=checkList] {Fire Function Repeatedly};
\node (held) [decision, below of=fireFunc] {Input Still Held?};

% Arrows
\draw [arrow] (start) -- (detect);

\draw [arrow] (detect) -- node[right] {Yes} (processEvent);
\draw [arrow] (detect.west) -- ++(-3,0) |- node[left] {No} (detect);

\draw [arrow] (processEvent) -- node[above] {Yes} (handleEvent);
\draw [arrow] (processEvent) -- node[right] {No} (checkList);

\draw [arrow] (handleEvent) |- (detect);

\draw [arrow] (checkList) -- (fireFunc);
\draw [arrow] (fireFunc) -- (held);

\draw [arrow] (held) -- node[right] {Yes} (fireFunc);
\draw [arrow] (held.west) -- ++(-3,0) |- node[left] {No} (detect);

\end{tikzpicture}


\subsubsection{Dependencies}

\subsubsection{Usage-example}

\subsubsection{Known issues and bottle-necks}

\subsubsection{Performance}